\documentclass[12pt]{article}
\usepackage{graphicx}
\usepackage[T1]{fontenc}
\usepackage{tocloft}
\usepackage{todonotes}
\usepackage{subcaption}
\usepackage{caption}
\usepackage{hyperref}
\usepackage{booktabs}
\usepackage{listings}
\usepackage{wrapfig}
\usepackage{pdfpages}
\usepackage{pdflscape}
\usepackage{scrlayer}
\usepackage{scrhack}

\begin{document}
\begin{titlepage}
    \centering
    {\scshape\large AY 2020/2021 \par}
    \vfill
    \includegraphics[width=100pt]{assets/logo-polimi-new}\par\vspace{1cm}
    {\scshape\LARGE Politecnico di Milano \par}
    \vspace{1.5cm}
    {\huge\bfseries RASD\@: Requirement Analysis
        and Specification Document \par}
    \vspace{2cm}
    {\Large {Alice Piemonti\quad Luca Pirovano\quad Nicolò Sonnino}\par}
    \vfill
    {\large Professor\par
        Matteo \textsc{Rossi}}
    \vfill
    {\large \today \par}
\end{titlepage}
\hypersetup{%
    pdfborder = {0 0 0}
}
\thispagestyle{plain}
\pagenumbering{gobble}
\mbox{}
\newpage
\pagenumbering{roman}
\tableofcontents
\newpage
\pagenumbering{arabic}
\section{Introduction}
CLup (Customers Line-up) is an easy-to-use application which intent is to help grocery shopping to face the big challenges that arise during this tough period.

In fact, due to the recent worldwide spread of SARS-CoV-2 (COVID-19), many countries are imposing lockdowns and strict regulations about social distancing such as: the closure of restaurants in the evening, limitations on public transports, curfews, etc.\\

In particular, supermarkets are imposed to restrict accesses to their stores, in order to avoid having too much people inside them, which stands at the basis of social distancing. Furthermore, the staggered access needed by stores lead to the same problem in the outside, such as long queues and crowds.

Grocery stores need an application which main intent is to manage in an efficient way customers arrival on the outside of the supermarket, and to count the access of people inside them. The advantage is to regulate the influx of people near the stores, according to the vigent strict rules, and to avoid long hours of waiting, saving people from line up.

\subsection{Purpose}
The aim of the product is to avoid gatherings outside and inside grocery stores,  improving the safety of the customers.

This is achieved through monitoring accesses to the buildings, managing time slots for visits and optimizing people flows inside the stores.\\

The application should provide two types of accesses for customers (as normal users) and stores' attendants (as special users).

Customers will have the possibility to line up in a virtual queue, so that they can wait from a close and safe building until their number is called. In a reasonable time, the application will inform the user when his number is about to be called, so that he can reach the store in the right time.\\

In addition, users can book a visit to the supermarket for a different time or day. Hence, the customer indicates the slot preferred and, if it is available, the application register the reservation; otherwise a list of alternatives are displayed, such as the possibility to book in a different slot, or to chose different stores available for that time/day.\\

The application can also register in advance the duration of the visit of a customer, and a list of categories that the customer intends to buy, so that the system can suggest the best slot and plan in a finer way the visits, in order to guarantee a correct distance through the aisles inside the store. In fact, the application will be able to balance the presence of people in all the areas of the supermarket, as well as the flows of visitors throughout the day.

Customers will be able to activate an additional functionality: the application will send a notification of available slots in a certain day/time range. Hence, the user will be facilitated in the search of an available slot, and this will guarantee that the user won't spend too much time managing the reservations.\\

Attendants will have the possibility to scan a QR code (generated by the user's application) at the entrance of the store with a specific functionality offered by the application itself; this helps the attendants to verify the correctness of customers' arrivals and, in the meanwhile, monitor the number of entrances.

In addition, attendants will have access to a proper area of the application that permits to behave as a proxy and hand out tickets on spot, in order to guarantee the access to those people who doesn't have access to the application (elderly people, people who don't have a smartphone or an Internet access).\\

The application will be operable freely, widely available and very intuitive, because the range of users (i.e. people who need to go to the grocery store) extends to the entire population.

The userbase is expected to be both people with an Internet access and ones without it, from young people to elderly, thanks to the possibility of attendants to act as a proxy.

\subsection{Scope}
The product shall be called CLup and will let users to plan their shopping session in two different ways:
\begin{itemize}
    \item {\textbf{ASAP}: the user will claim the first available ticket and receive an estimated queue time.}
    \item {\textbf{Reservation}: the user will choose a day/time slot from a list of available ones, in order to book his visit to the structure.}
\end{itemize}

Every customer can choose one of these modalities \textbf{remotely} via an official app or through a web browser, or \textbf{in presence} by asking to a staff member, who will act as an intermediate between the customer and the system.

When a customer makes a reservation, the system allows him to choose the duration of his visit and insert a list of possible purchases, in order to optimize his stay.\\

In addition to that, the user can change both the slot and the store relying on system's suggestions: the application will display the best recommendation in order to balance the number of people inside a store throughout the day, merging in the same time slots those people who will purchase different product, in order to guarantee a good distancing through the aisles as well. The user can also enable periodically notifications of available slots in a day/time range: this will allow the application to notify the user when a desirable slot (according to day/time preferences) is available.


\todo{aggiungere introduzione e nome sottosezione}
\begin{center}
    \begin{table}

        \begin{tabular}{|c|c|c|}
            \hline
            Phenomenon                                 & Who controls it? & Is shared? \\
            \hline
            The formation of queues                    & W                & N          \\
            Social distancing                          & W                & N          \\
            User registration                          & M                & Y          \\
            User login                                 & W                & Y          \\
            Check username and password                & M                & N          \\
            User retrieves a new number                & W                & Y          \\
            Generate a unique number                   & M                & N          \\
            Visualize queue number                     & W                & Y          \\
            Check time-out                             & M                & N          \\
            Notify number is going to be called        & M                & Y          \\
            Generate a QR code                         & M                & N          \\
            Visualize QR code                          & M                & Y          \\
            Attendant scans QR code                    & W                & Y          \\
            Check QR code validity                     & M                & N          \\
            Customer retrieves a number on spot        & W                & Y          \\
            User selects day/time of slot              & W                & Y          \\
            Check availability of slot                 & M                & N          \\
            User adds visit duration                   & W                & Y          \\
            Check visit duration validity              & M                & N          \\
            User indicates list if items               & W                & Y          \\
            Confirm slot has been reserved             & M                & Y          \\
            Cross data                                 & M                & N          \\
            Make recommendation of alternative slots   & M                & N          \\
            Visualize suggestions of alternative slots & M                & Y          \\
            Check customers balance                    & M                & N          \\
            Accept receipt of notifications            & W                & Y          \\
            Search periodically for available slots    & M                & N          \\
            Notify for available slots                 & M                & Y          \\

            \hline
        \end{tabular}
        \caption{phenomena table}
    \end{table}
\end{center}

\subsection{Goals}
The main objectives of our system are the following:
\begin{itemize}
    \item {\textbf{G1: Allow users to retrieve a unique queue number}\\This is the main feature of the application, through which customers are invited not staying outside the structure. Through an appropriate estimation of each customer's permanence time, the user receives an estimated queue time after which he can go to the store. The number is guaranteed to be unique.}
    \item {\textbf{G2: Allow users to generate a QR Code}\\The QR Code is needed in order to let store attendants to monitor entrances, scanning the customer's code upon entering.}
    \item {\textbf{G3: Allow shops to offer this service}\\This goal can be split into the following:
          \begin{itemize}
              \item {\textbf{Allow shops to offer the remote queue function}\\Each shop can register to the service and let customers to Line-Up from home. The application will track the state of the queue.}
              \item {\textbf{Allow shops to generate tickets on the spot}\\If someone does not have access to the required technology, they can still take advantage of the system by getting their queue number directly at the store, where an attendant will act as an intermediate.}
          \end{itemize}
          }
    \item {\textbf{G4: Allow users to "book a visit"}\\Users can book a slot providing the expected duration of the visit. Moreover, users can provide a list which describes the categories of products they intend to buy. Thank to this, visits are optimized and those customers who belong to the same slot, are supposed to buy different types of product, so that they stay in different areas of the store.}
    \item {\textbf{G5: Let the system to infer a visit duration}\\For long term users, the system provides an estimation of visit duration, relying on their previous shopping sessions data.}
    \item {\textbf{G6: Allow users to receive a suggestion of alternative slots}\\During the reservation procedure users will receive a list of alternatives. This help to balance out the number of people in the store during the day. The suggestions could be both on the same shop chain and on different chains.}
    \item {\textbf{G7: Allow users to receive notifications on available slots}\\Users could enable notifications in order to receive periodic announcements of available slots in a preferred day/time range.}
\end{itemize}
\subsection{Definitions, Acronyms, Abbreviations}
\subsubsection{Definitions}
\begin{itemize}
    \item {\textbf{Customers:} the people whom this service is directed. They can belong to any age and sex. Their main purpose is to request a ticket to schedule their line up at the shop.}
    \item {\textbf{User:} it is a customer who uses the application. }
    \item {\textbf{Store:} it is the provider of the goods that customers want to buy. It is obliged to limit the number of people entering in its building through a line up method, in order to respect the local laws for COVID-19 pandemic.}
    \item {\textbf{Attendant:} it is a store's employee. Its aim is both to manage entrances and to help customers. For this reasons attendants can scan QR codes from customers' application and also release tickets on spot acting as intermediate for customers.}
    \item {\textbf{Queue number:} it is the ID assigned to each user. It's composed of alphanumerical elements, and it's needed to access the desired shop. It can be retrieved online or on the spot.}
    \item {\textbf{Visit:} it refers to the customers entering the shop, and also to their staying time. It is associated to both a certain store and a day/time range.}
    \item {\textbf{Slot:} it is a day/time range. It can be reserved by a limit number of users in order to guarantee a maximum number of people who are inside a store in every time of the day.}
    \item {\textbf{QR Code:} it's a graphic representation of a string, which can be easily read through barcode readers. It contains the visual representation of the queue number, in order to let attendants make a customer "check-in".}
    \item {\textbf{Notification:} it's an alert that a certain event occurred. This alert can be a "Push Notification" on the smartphone, an SMS, an email and so on.}
    \item {\textbf{Push notification:} it's an automated message sent by an application to a user when the application is not running.}
\end{itemize}
\subsubsection{Acronyms}
\begin{itemize}
    \item {\textbf{ASAP:} As Soon As Possible. It refers to the possibility of getting an appointment on the first available slot.}
    \item {\textbf{S2B:} Software to Be, it's the one designed in this document and not yet implemented.}
    \item {\textbf{API:} Application Programming Interface, it indicates on demand procedure which supply a specific task.}
\end{itemize}
\subsubsection{Abbreviations}
\begin{itemize}
    \item {\textbf{ID:} identifier. It's a generally unique sequence of numbers or letters in order to unambiguously identify an entity.}
    \item {\textbf{Gn:} goal number n}
\end{itemize}
\subsection{Revision History}
\subsection{Reference Documents}
\begin{itemize}
    \item {Specification document: "R\&DD Assignment A.Y. 2020-2021"}
    \item {Alloy official documentation: \href{https://alloytools.org/documentation.html}{https://alloytools.org/documentation.html}}
    \item {Paper: "Jackson and Zave: the world and the machine"}
\end{itemize}
\subsection{Document Structure}
\begin{itemize}
    \item {\textbf{Section 1: Introduction}\\This section offers a brief description of the problem and required functionalities. \\It also contains the list of definitions, acronyms and abbreviations that could be found in this document. \\Finally, there are changelog of the document, containing the revisions list and their content, and document structure, which describes the main purposes of the sections of this document.}
    \item {\textbf{Section 2: Overall Description}\\This section offers a summary description about interactions between the S2B and the external interfaces.}
    \item {\textbf{Section 3: Specific Requirements}}
    \item {\textbf{Section 4: Formal Analysis through Alloy}}
\end{itemize}
\newpage
\section{Overall Description}
\subsection{Perspective}
CLup is a complete system that offers the functionalities described in the specific \textit{Product Functions} section.

It does not lean on a previously existent system and does not offer integrations with other existing systems.

CLup provides all the module that are necessary for its execution. However, it must exploit external interfaces to accomplish its requirements.

\subsubsection{User Interfaces}
The system should interfaces with users through devices which must be connected to Internet.

Everyone that needs to use this service would connect to the service through a Web Interface (through an existent domain, like \textit{www.clup.com}) or through a mobile application that can be installed on smartphones (both IOS and Android).

\subsubsection{Software Interfaces}
\subsubsection{Hardware Interfaces}
The system "as it is" does not provide any hardware interface. However, local shops could use the QR Code given by the application on their local systems in order to monitor the number of customers who have entered the shop.
\subsubsection{Hardware Constraints}
Each person interested in using the system described in this document to line up for a shop should have a device connected to Internet, which could be a smartphone, tablet or personal computer.

The basic requirement is to have an Internet browser installed on the used device.

Shops, instead, should have some devices connected to the service (through web interface or application) in order to monitor queue length, entrances and, if needed, close the booking session.

\subsection{Product Functions}

\subsection{Actors}
\subsubsection{User}
It's a customer, i.e., a normal person who intend to go to the grocery shop. It is able to use the smartphone application or surf on the website of the application. It uses the service in order to access the stores, according to the regulations. It takes advantage in using the application, since, when it want to access a store, it retrieves the unique number and can wait from a close and safe building.

\subsubsection{Attendant}
It's a worker of a specific store. Its job is to help customers and monitor the entrances of people. It uses the application for scanning the QR code of users at the entrance of the store. In addition, it can enter in a dedicated section of the application in order to act as an intermediate for those customers who don't have an Internet access.

\subsubsection{Store}
It's a shop. For the application's purpose only, we identify stores as grocery shops. It adopts the service in order to manage the flows of people inside and outside the building, as well as observing the laws on social distancing.

\subsection{Assumption, Dependencies, Constraints}

\newpage
\section{Specific Requirements}
\subsection{Interface Requirements}
\subsection{Functional Requirements}
\subsubsection{User Scenarios}
\textbf{Scenario 1}\\

\newpage
\section{Formal Analysis}
\newpage
\section{Effort spent}
\begin{tabular}{ | l || c | r | l | c |}
    \hline
    Student        & Time for S.1 & Time for S.2 & Time for S.3 & Time for S.4 \\ \hline
    Alice Piemonti & 3h           &              &              &              \\ \hline
    Luca Pirovano  & 2h25         & 0h40         &              &              \\ \hline
    Nicolò Sonnino &              &              &              &              \\
    \hline
\end{tabular}
\end{document}